\section{Resultados por fase del pipeline NeuroMIND}

Esta sección incorpora una primera lectura técnica de los artefactos generados por \path{EEG_Julia/NeuroMIND}. La ejecución revisada corresponde a \texttt{sub-M05}, sesión \texttt{T2}, tarea \texttt{eyesclosed} (condición EC), con \SI{500}{Hz}, 31 canales y \SI{100.36}{s} de señal. Debe leerse como validación operativa de la cadena de análisis y no como evidencia clínica concluyente: no hay comparación grupal ni contraste EO/EC en esta salida. La actualización actual ya incorpora inferencia por surrogates de fase aleatorizada y corrección FDR para wPLI.

\begin{decisionmetodologica}
Las matrices actuales de NeuroMIND están calculadas en espacio de sensores con \texttt{use\_csd=false}. Por tanto, las figuras wPLI de esta sección se interpretan como conectividad funcional sensor--sensor preliminar. Los surrogates/FDR ya permiten distinguir patrones observados de evidencia inferencial dentro de esta ejecución, pero la lectura definitiva sigue requiriendo relanzar la configuración final con CSD y comparación entre sujetos/condiciones.
\end{decisionmetodologica}

\subsection{Resumen fase a fase de la ejecución}

La \cref{tab:results-phase-dashboard-reading} resume la evidencia generada por la ejecución y la conecta con la organización metodológica de la \cref{tab:pipeline-operational-contracts}. La intención no es duplicar todos los CSV, sino fijar qué dato sostiene cada interpretación del informe y qué parte debe revisarse en el dashboard cuando aparezca una anomalía.

\begin{fullwidth}
\begin{table}[H]
\centering
\scriptsize
\caption{Lectura fase a fase de resultados para \texttt{sub-M05}/\texttt{ses-T2}/\texttt{eyesclosed}. Las métricas resumen proceden de los artefactos exportados por NeuroMIND y se interpretan como control técnico de una ejecución individual.}
\label{tab:results-phase-dashboard-reading}
\begin{adjustbox}{max width=0.99\fulltextwidth}
\begin{tabular}{@{}p{0.12\fulltextwidth}p{0.22\fulltextwidth}p{0.23\fulltextwidth}p{0.24\fulltextwidth}p{0.12\fulltextwidth}@{}}
\toprule
\textbf{Fase} & \textbf{Artefacto revisado} & \textbf{Resultado observado} & \textbf{Interpretación técnica prudente} & \textbf{Panel} \\
\midrule
IO/QC & \texttt{overview.csv}, \texttt{qc\_summary.csv} & 31 canales, \SI{500}{Hz}, \SI{100.36}{s}; \texttt{Fp2} marcado como malo por RMS & Ejecución suficiente para auditar pipeline; las lecturas frontales deben contrastarse con QC & 0--1 \\
\addlinespace
Filtrado & vista raw/filtrada, configuración & Señal más centrada tras filtrado; persisten transitorios locales & El filtrado estabiliza la señal, pero no sustituye rechazo de artefactos ni revisión visual & 2--4 \\
\addlinespace
ICA & resumen ICA, features IC & 31 componentes; 0 rechazados; 26 brain, 3 line noise, 2 jump & La clasificación está disponible, pero la política de rechazo fue conservadora en esta corrida & 5 \\
\addlinespace
Segmentación & resumen de épocas y rechazados & 98/100 épocas válidas; 2 rechazos por amplitud; cobertura media 99.7\% & Base temporal razonable para PSD/wPLI, con trazabilidad de los segmentos excluidos & 6--7 \\
\addlinespace
Espectral & PSD, potencia por banda & Delta 38.6\%, alfa 25.51\%, theta 14.36\%; alfa posterior destacada & La estructura espectral orienta qué bandas revisar, pero no demuestra conectividad & 8 \\
\addlinespace
wPLI/red & matrices, aristas, métricas nodales & Alfa tiene wPLI medio 0.3317 y densidad exploratoria 0.8559 & Patrón descriptivo sensor--sensor; requiere CSD y contraste inferencial para lectura científica & 9 \\
\addlinespace
Surrogates/FDR & p/q-valores, máscaras, conexiones & 200 surrogates, BH-FDR, 0/465 aristas significativas por banda & La evidencia inferencial de esta ejecución individual es negativa & 10 \\
\bottomrule
\end{tabular}
\end{adjustbox}
\end{table}
\end{fullwidth}

\begin{controlqc}
La sección de resultados queda organizada como una cadena de custodia: cada figura se lee contra su tabla o resumen exportado. Si una fase falla, las fases posteriores no se eliminan automáticamente, pero su interpretación cambia de ``resultado neurofisiológico'' a ``evidencia de depuración del pipeline''.
\end{controlqc}

\subsection{Fases 0--1: IO/BIDS, QC inicial y filtrado}

La salida de control de calidad identifica un canal marcado como problemático (\texttt{Fp2}, $z_{\mathrm{RMS}}=4.63$), mientras que el resto de canales quedan dentro de los umbrales automáticos. La segmentación retiene 98 de 100 épocas de \SI{1}{s}, con dos rechazos por amplitud y una cobertura media del 99.7\%. Estos indicadores sugieren que la ejecución es útil para revisar el funcionamiento del pipeline, aunque el canal frontal señalado debe considerarse al interpretar potencia lenta y conectividad que lo involucre.

\begin{fullwidth}
\begin{figure}[H]
\centering
\includegraphics[width=0.52\fulltextwidth]{results/neuromind_m05_t2_ec_senal_raw_filtrada_crop.png}
\caption{Vista previa de señal en NeuroMIND para \texttt{M05/T2/EC}. La curva gris muestra fragmentos raw y la azul la señal filtrada. La figura permite comprobar visualmente que la cadena notch/bandreject/high-pass/low-pass reduce deriva y componentes fuera de banda, pero también muestra que algunos canales frontotemporales conservan transitorios relevantes para QC.}
\label{fig:results-neuromind-signal-preview}
\end{figure}
\end{fullwidth}

\textbf{Resultado observado.} La señal filtrada queda más centrada y con menor deriva lenta que la raw en los canales de ejemplo, pero la variabilidad intercanal no desaparece por completo.

\textbf{Interpretación técnica.} La figura es coherente con un preprocesamiento diseñado para estabilizar la entrada de ICA, segmentación y análisis espectral, no para eliminar toda estructura temporal. Los fragmentos con amplitud o morfología atípica justifican mantener controles por canal y por época antes de propagar datos a PSD y wPLI.

\textbf{Limitación.} La vista previa no resume todos los canales ni todas las épocas; debe complementarse con tablas QC, histogramas de amplitud y trazabilidad de épocas rechazadas.

\subsection{Fases 2--5: ICA, segmentación, baseline y rechazo de artefactos}

La etapa ICA exporta 31 componentes, conserva el 100\% de la varianza y no rechaza componentes en esta ejecución. La clasificación automática etiqueta 26 componentes como compatibles con actividad cerebral, 3 con ruido de línea y 2 con saltos/transitorios. Este resultado debe entenderse como una clasificación disponible para inspección, no como una limpieza agresiva: la reconstrucción conserva todos los componentes y traslada parte del control a la segmentación y al rechazo por época.

\begin{table}[H]
\centering
\caption{Resumen operativo de ICA y segmentación en la ejecución revisada. La ausencia de componentes rechazados indica una política conservadora de limpieza ICA.}
\label{tab:results-ica-segmentation}
\begin{tabular}{ll}
\toprule
\textbf{Indicador} & \textbf{Valor} \\
\midrule
Componentes ICA calculados & 31 \\
Componentes rechazados & 0 \\
Clasificación IC & 26 brain, 3 line noise, 2 jump \\
Épocas totales / válidas & 100 / 98 \\
Longitud de época & \SI{1}{s} \\
Rechazos por amplitud / gradiente & 2 / 0 \\
Canales usados para rechazo & 30 de 31 \\
Cobertura media / mínima & 99.7\% / 94.0\% \\
\bottomrule
\end{tabular}
\end{table}

\begin{fullwidth}
\begin{figure}[H]
\centering
\begin{subfigure}[t]{0.23\fulltextwidth}
\centering
\includegraphics[width=\linewidth]{results/ica/neuromind_m05_t2_ec_ica_ic05_jump.png}
\caption{IC05: jump.}
\end{subfigure}
\hfill
\begin{subfigure}[t]{0.23\fulltextwidth}
\centering
\includegraphics[width=\linewidth]{results/ica/neuromind_m05_t2_ec_ica_ic13_jump.png}
\caption{IC13: jump.}
\end{subfigure}
\hfill
\begin{subfigure}[t]{0.23\fulltextwidth}
\centering
\includegraphics[width=\linewidth]{results/ica/neuromind_m05_t2_ec_ica_ic14_line_noise.png}
\caption{IC14: línea.}
\end{subfigure}
\hfill
\begin{subfigure}[t]{0.23\fulltextwidth}
\centering
\includegraphics[width=\linewidth]{results/ica/neuromind_m05_t2_ec_ica_ic21_line_noise.png}
\caption{IC21: línea.}
\end{subfigure}
\caption{Topoplots ICA exportados por NeuroMIND para componentes con puntuaciones automáticas altas en la ejecución revisada. El círculo representa la cabeza vista desde arriba, con la nariz en posición superior; los puntos son electrodos y el color indica el peso relativo de cada IC sobre el cuero cabelludo. IC05 e IC13 quedan clasificados como transitorios/saltos, mientras que IC14 e IC21 concentran evidencia compatible con ruido de línea. En esta corrida no se rechazaron componentes, por lo que la figura funciona como auditoría de clasificación y no como demostración de limpieza aplicada.}
\label{fig:results-ica-topomaps-selected}
\end{figure}
\end{fullwidth}

\textbf{Cómo leer estos mapas.} A diferencia de las matrices wPLI, estos paneles no son heatmaps de canal por canal. Cada topoplot resume una columna de la matriz de mezcla ICA: muestra con qué polaridad e intensidad se proyecta un componente independiente sobre los electrodos. Las zonas rojas y azules no significan activación o desactivación cerebral en sentido anatómico; indican pesos positivos o negativos de una fuente separada por ICA. Por ello, la lectura útil es geométrica y diagnóstica: se busca decidir si el componente parece una fuente neural plausible, un parpadeo, un movimiento ocular, músculo, ruido de línea o un salto de electrodo.

\begin{table}[H]
\centering
\scriptsize
\caption{Comentario individual de los topoplots ICA seleccionados. La etiqueta procede de la clasificación automática de NeuroMIND y la interpretación combina la geometría visible del mapa con la familia de score asignada.}
\label{tab:results-ica-topomap-comments}
\begin{adjustbox}{max width=0.98\textwidth}
\begin{tabular}{@{}p{0.10\textwidth}p{0.14\textwidth}p{0.35\textwidth}p{0.31\textwidth}@{}}
\toprule
\textbf{IC} & \textbf{Etiqueta} & \textbf{Firma topográfica observada} & \textbf{Lectura para esta ejecución} \\
\midrule
IC05 & \texttt{jump} & Foco rojo muy localizado en el borde temporal derecho, con campo opuesto azul amplio y suave hacia el hemisferio izquierdo. No muestra el patrón frontal bilateral típico de parpadeo ni un dipolo posterior alfa limpio. & Compatible con una contribución local/periférica o transitoria, posiblemente ligada a sensor o movimiento breve. La etiqueta \texttt{jump} debe confirmarse con la traza temporal del IC; al no haberse rechazado, actúa como componente a vigilar en auditoría. \\
\addlinespace
IC13 & \texttt{jump} & Mancha azul pequeña e intensa en el borde inferior derecho/posterior, con el resto del cuero cabelludo prácticamente neutro. La concentración extrema en una zona periférica pesa más que una distribución neural suave. & Patrón típico de componente localizado, poco compatible con fuente cerebral distribuida. La lectura prudente es transitorio o irregularidad local; conviene revisar si coincide con picos de amplitud en épocas rechazadas. \\
\addlinespace
IC14 & \texttt{line\_noise} & Patrón irregular de baja amplitud global, con foco azul temporal derecho y focos rojizos débiles cercanos. La geometría no basta por sí sola para diagnosticar línea, porque el ruido eléctrico se reconoce sobre todo por el espectro. & La etiqueta depende principalmente del score espectral de línea. El topoplot no sugiere un componente cerebral canónico; debe inspeccionarse el pico en 50\,Hz o armónicos antes de decidir eliminación en una corrida final. \\
\addlinespace
IC21 & \texttt{line\_noise} & Campo rojizo muy extendido con máximo frontal superior y pequeñas zonas azuladas centrales/laterales. No reproduce el dipolo frontolateral opuesto típico de HEOG ni la simetría posterior esperada de alfa EC. & Candidato a contaminación global o componente no fisiológico con evidencia de línea. Como en IC14, la decisión no debe descansar solo en el mapa: la PSD del IC y el score \texttt{line\_score} son los criterios determinantes. \\
\bottomrule
\end{tabular}
\end{adjustbox}
\end{table}

\begin{controlqc}
\small
La interpretación visual de ICA debe evitar dos errores: tratar el topoplot como una imagen directa de actividad cerebral y rechazar componentes únicamente por color intenso. Un componente cerebral real puede tener pesos fuertes si forma un dipolo suave y fisiológicamente plausible; un artefacto puede tener pesos modestos pero una firma espectral o temporal inequívoca. En NeuroMIND, la decisión final debe cruzar topografía, PSD, activación temporal y scores automáticos antes de modificar la lista de componentes rechazados.
\end{controlqc}

\textbf{Resultado observado.} Los segmentos rechazados se asocian a \texttt{C4} y \texttt{FT9}, ambos con un 1\% de épocas marcadas por amplitud; \texttt{Fp2} queda excluido del criterio de rechazo por estar marcado como canal problemático.

\textbf{Interpretación técnica.} La cadena es coherente con una estrategia de protección de datos: ICA no elimina componentes automáticamente y la segmentación actúa como control posterior. Esto es útil para auditoría, aunque deja abierta la necesidad de inspección manual o criterios más explícitos antes de una ejecución definitiva.

\textbf{Limitación.} La clasificación ICA automática no debe tratarse como veredicto fisiológico. Para el informe final conviene conservar topografías, espectros y cursos temporales de componentes rechazados si en futuras ejecuciones se activa una política de exclusión.

\subsection{Fase 6: análisis espectral y potencia por banda}

El resumen espectral usa ventana Hamming, $N_{\mathrm{fft}}=1024$, $\Delta f\approx\SI{0.488}{Hz}$ y referencia promedio. En la ejecución revisada, la potencia media agregada está dominada por delta (38.6\% del total) y alfa (25.51\%), seguida de theta (14.36\%). La potencia alfa más alta aparece en regiones occipitales/parieto-occipitales, un patrón técnicamente compatible con reposo en ojos cerrados, aunque no debe transformarse aquí en conclusión neurofisiológica individual.

\begin{fullwidth}
\begin{figure}[H]
\centering
\includegraphics[width=0.94\fulltextwidth]{results/neuromind_m05_t2_ec_psd_canales.png}
\caption{Densidad espectral de potencia por canal. Cada panel muestra la distribución frecuencial estimada para un electrodo, lo que permite inspeccionar picos dominantes, diferencias regionales y canales con energía anómala antes de resumir la señal por bandas.}
\label{fig:results-neuromind-psd}
\end{figure}
\end{fullwidth}

\begin{fullwidth}
\begin{figure}[H]
\centering
\includegraphics[width=0.84\fulltextwidth]{results/neuromind_m05_t2_ec_potencia_bandas.png}
\caption{Potencia por banda y canal. La visualización condensa la PSD en bandas funcionales, mostrando una distribución alfa marcada en electrodos posteriores y elevaciones lentas localizadas que deben revisarse junto con QC, especialmente cuando afectan a canales frontales o a canales marcados como potencialmente problemáticos.}
\label{fig:results-neuromind-band-power}
\end{figure}
\end{fullwidth}

\textbf{Resultado observado.} La potencia alfa alcanza máximos posteriores, con valores altos en \texttt{O2}, \texttt{P8}, \texttt{O1} y \texttt{Oz}; delta presenta valores elevados en \texttt{Fp2} y \texttt{C4}.

\textbf{Interpretación técnica.} La distribución alfa posterior favorece que la banda alfa sea una candidata natural para revisar conectividad funcional en EC. Las elevaciones delta localizadas, en cambio, deben tratarse con cautela porque pueden reflejar actividad lenta real, contaminación residual, referencia, movimiento o efectos de canal.

\textbf{Limitación.} La potencia por banda no es una prueba de conectividad: solo informa sobre energía local por canal. Para inferir acoplamiento funcional se requiere analizar relaciones entre pares y controlar multiplicidad.

\subsection{Fase 7: matrices wPLI observadas}

NeuroMIND exporta matrices wPLI por banda con 465 aristas posibles para 31 canales. En esta ejecución, alfa presenta la media wPLI más alta (0.3317), el máximo más alto (0.868) y la mayor densidad por encima del umbral exploratorio 0.1 (0.8559). Theta muestra una densidad intermedia-alta (0.6495), pero con máximos más bajos (0.4347). Estos valores describen estructura observada de sincronización de fase; su lectura inferencial se evalúa después con surrogates y FDR.

\begin{fullwidth}
\begin{figure}[H]
\centering
\begin{subfigure}[t]{0.48\fulltextwidth}
\centering
\includegraphics[width=\linewidth]{results/neuromind_m05_t2_ec_wpli_alpha.png}
\caption{wPLI alfa.}
\end{subfigure}
\hfill
\begin{subfigure}[t]{0.48\fulltextwidth}
\centering
\includegraphics[width=\linewidth]{results/neuromind_m05_t2_ec_wpli_theta.png}
\caption{wPLI theta.}
\end{subfigure}
\caption{Matrices wPLI por banda para \texttt{M05/T2/EC}. La banda alfa muestra valores altos y estructura posterior más marcada; theta presenta conectividad más moderada y distribuida. La diagonal nula y la simetría reflejan que wPLI se calcula como métrica no dirigida entre pares de canales.}
\label{fig:results-neuromind-wpli-alpha-theta}
\end{figure}
\end{fullwidth}

\begin{fullwidth}
\begin{figure}[H]
\centering
\begin{subfigure}[t]{0.31\fulltextwidth}
\centering
\includegraphics[width=\linewidth]{results/wpli/neuromind_m05_t2_ec_wpli_delta.png}
\caption{Delta.}
\end{subfigure}
\hfill
\begin{subfigure}[t]{0.31\fulltextwidth}
\centering
\includegraphics[width=\linewidth]{results/wpli/neuromind_m05_t2_ec_wpli_beta_low.png}
\caption{Beta baja.}
\end{subfigure}
\hfill
\begin{subfigure}[t]{0.31\fulltextwidth}
\centering
\includegraphics[width=\linewidth]{results/wpli/neuromind_m05_t2_ec_wpli_beta_mid.png}
\caption{Beta media.}
\end{subfigure}

\vspace{0.35em}

\begin{subfigure}[t]{0.31\fulltextwidth}
\centering
\includegraphics[width=\linewidth]{results/wpli/neuromind_m05_t2_ec_wpli_beta_high.png}
\caption{Beta alta.}
\end{subfigure}
\hfill
\begin{subfigure}[t]{0.31\fulltextwidth}
\centering
\includegraphics[width=\linewidth]{results/wpli/neuromind_m05_t2_ec_wpli_gamma.png}
\caption{Gamma.}
\end{subfigure}
\caption{Atlas complementario de matrices wPLI exportadas por el pipeline para bandas no mostradas en la figura principal. El panel permite comparar visualmente la densidad y escala relativa de conectividad observada entre bandas, pero debe leerse junto con la tabla surrogate/FDR: ninguna banda retiene aristas significativas tras corrección de multiplicidad en esta ejecución individual.}
\label{fig:results-neuromind-wpli-band-atlas}
\end{figure}
\end{fullwidth}

\begin{table}[H]
\centering
\caption{Aristas wPLI superiores en bandas seleccionadas y contraste surrogate asociado. Los valores altos observados no implican significación cuando el p-valor empírico y el q-valor FDR permanecen altos.}
\label{tab:results-top-wpli-edges}
\begin{tabular}{llllrrr}
\toprule
Banda & Canal A & Canal B & Lectura técnica & wPLI & $p_{\mathrm{emp}}$ & $q_{\mathrm{FDR}}$ \\
\midrule
ALPHA & Oz & P4 & posterior/parietal & 0.8680 & 0.755 & 0.8015 \\
ALPHA & O1 & P4 & occipito-parietal & 0.8675 & 0.755 & 0.8015 \\
ALPHA & O1 & O2 & interoccipital & 0.8612 & 0.610 & 0.7985 \\
THETA & Fz & Fp2 & frontal, revisar QC & 0.4347 & 0.420 & 0.9221 \\
THETA & FC1 & Fp2 & frontocentral, revisar QC & 0.4015 & 0.580 & 0.9221 \\
\bottomrule
\end{tabular}
\end{table}

\textbf{Resultado observado.} Las aristas alfa superiores se concentran en pares occipitales/parietales, mientras que algunas aristas theta superiores implican \texttt{Fp2}, canal marcado como malo por QC.

\textbf{Interpretación técnica.} La matriz alfa es coherente con la sección teórica: wPLI capta consistencia de desfase no nulo entre señales analíticas y puede resaltar relaciones posteriores en reposo EC. La implicación de \texttt{Fp2} en theta ilustra por qué conectividad y QC deben interpretarse juntos: un valor wPLI alto no garantiza validez fisiológica si una señal de entrada está contaminada.

\textbf{Limitación.} Aunque las matrices muestran patrones observados, el contraste surrogate de esta ejecución no respalda aristas significativas tras FDR. Por tanto, las aristas destacadas son útiles para depurar e interpretar el pipeline, no para formular conclusiones clínicas.

\subsubsection*{Lectura de red derivada}

Además de las matrices, NeuroMIND exporta métricas nodales calculadas sobre la red wPLI umbralizada de forma exploratoria. En la ejecución revisada, los mayores valores de fuerza normalizada aparecen en \texttt{C4}, \texttt{FT10}, \texttt{CP6}, \texttt{T8}, \texttt{P8} y \texttt{FC2}. Esta lectura ayuda a localizar qué nodos concentran peso de conectividad en el grafo descriptivo, pero no cambia el resultado inferencial posterior.

\begin{table}[H]
\centering
\caption{Nodos con mayor fuerza normalizada en la red wPLI exploratoria. Estos valores describen centralidad dentro de la matriz observada, no significación estadística.}
\label{tab:results-network-nodes}
\begin{tabular}{lrrr}
\toprule
Canal & Fuerza & Grado & Fuerza norm. \\
\midrule
C4 & 6.0767 & 20 & 1.0000 \\
FT10 & 5.3504 & 20 & 0.8805 \\
CP6 & 5.2567 & 17 & 0.8651 \\
T8 & 5.1256 & 19 & 0.8435 \\
P8 & 5.0257 & 16 & 0.8270 \\
FC2 & 4.9361 & 23 & 0.8123 \\
\bottomrule
\end{tabular}
\end{table}

\textbf{Resultado observado.} El máximo nodal aparece en \texttt{C4}, un canal que también aparece en el resumen de épocas rechazadas por amplitud.

\textbf{Interpretación técnica.} La coincidencia entre centralidad alta y trazas QC refuerza una idea metodológica importante: las métricas de red deben leerse con la misma cautela que las matrices de conectividad. Un nodo central puede reflejar una organización funcional relevante, pero también sensibilidad a artefactos residuales, referencia o umbralización.

\textbf{Limitación.} La tabla resume una red exploratoria y no una red estadísticamente validada. Para comparaciones entre grupos o sesiones, las métricas de red deben recalcularse con una convención fija de umbral, CSD y control de multiplicidad.

\subsection{Fase 8: inferencia con surrogates y FDR}

La actualización de NeuroMIND añade una prueba nula por aleatorización de fase: para cada banda y par de canales se generan 200 surrogates, se recalcula wPLI y se obtiene un p-valor empírico. Después se aplica corrección Benjamini--Hochberg con $\alpha=0.05$ sobre las 465 aristas de cada banda. La ejecución usa semilla 42 y exporta matrices de p-valores, q-valores, máscaras significativas y estadísticas nulas por banda.

\begin{fullwidth}
\begin{figure}[H]
\centering
\resizebox{0.94\fulltextwidth}{!}{%
\begin{tikzpicture}
\begin{axis}[
  width=13.8cm,
  height=5.4cm,
  ybar,
  bar width=5pt,
  ymin=0,
  ymax=0.70,
  symbolic x coords={Delta,Theta,Alpha,Beta low,Beta mid,Beta high,Gamma},
  xtick=data,
  x tick label style={rotate=35,anchor=east,font=\scriptsize},
  ylabel={valor medio},
  legend style={at={(0.5,1.08)},anchor=south,legend columns=3,draw=none},
  ymajorgrids=true,
  grid style={NSLine!45},
  tick label style={font=\scriptsize},
  label style={font=\small},
]
\addplot[fill=NSPetrol!70,draw=NSPetrol!90!black] coordinates {
  (Delta,0.1268) (Theta,0.1491) (Alpha,0.3317) (Beta low,0.1407) (Beta mid,0.1155) (Beta high,0.0864) (Gamma,0.0948)
};
\addplot[fill=NSOrange!65,draw=NSOrange!90!black] coordinates {
  (Delta,0.1401) (Theta,0.1746) (Alpha,0.3513) (Beta low,0.1681) (Beta mid,0.1481) (Beta high,0.0933) (Gamma,0.0989)
};
\addplot[fill=NSPurple!55,draw=NSPurple!85!black] coordinates {
  (Delta,0.5622) (Theta,0.6494) (Alpha,0.6171) (Beta low,0.6020) (Beta mid,0.6330) (Beta high,0.5764) (Gamma,0.5595)
};
\legend{wPLI observado medio, media nula surrogate, p medio}
\end{axis}
\end{tikzpicture}}
\caption{Resumen inferencial por banda para \texttt{M05/T2/EC}. Las medias wPLI observadas quedan próximas o por debajo de las medias nulas surrogate, y los p-valores medios son altos en todas las bandas. Esto explica que la corrección FDR no retenga aristas significativas.}
\label{fig:results-surrogate-band-summary}
\end{figure}
\end{fullwidth}

\begin{table}[H]
\centering
\caption{Control surrogate/FDR por banda. En todas las bandas, el umbral FDR efectivo es 0 y no queda ninguna arista significativa.}
\label{tab:results-surrogate-summary}
\begin{tabular}{lrrrrrr}
\toprule
Banda & wPLI medio & media nula & $p$ medio & $p_{\min}$ & $q_{\min}$ & aristas sig. \\
\midrule
Delta & 0.1268 & 0.1401 & 0.5622 & 0.075 & 0.9746 & 0/465 \\
Theta & 0.1491 & 0.1746 & 0.6494 & 0.280 & 0.9221 & 0/465 \\
Alpha & 0.3317 & 0.3513 & 0.6171 & 0.340 & 0.7985 & 0/465 \\
Beta baja & 0.1407 & 0.1681 & 0.6020 & 0.170 & 0.9406 & 0/465 \\
Beta media & 0.1155 & 0.1481 & 0.6330 & 0.295 & 0.8868 & 0/465 \\
Beta alta & 0.0864 & 0.0933 & 0.5764 & 0.395 & 0.7907 & 0/465 \\
Gamma & 0.0948 & 0.0989 & 0.5595 & 0.270 & 0.7888 & 0/465 \\
\bottomrule
\end{tabular}
\end{table}

\textbf{Resultado observado.} La tabla de conexiones significativas exportada por NeuroMIND está vacía: \texttt{n\_sig\_total = 0}. Incluso las aristas con wPLI observado alto en alfa tienen p-valores empíricos elevados, lo que indica que no se separan de forma suficiente de la distribución nula generada por aleatorización de fase.

\textbf{Interpretación técnica.} El resultado no invalida la matriz wPLI observada; la reubica. La matriz sirve para estudiar estructura descriptiva y revisar si el pipeline produce patrones esperables, mientras que la inferencia surrogate/FDR pregunta si esos patrones superan una nula conservadora para esta ejecución individual. En esta corrida, la respuesta inferencial es negativa.

\textbf{Limitación.} El contraste se aplica a un único sujeto/sesión/condición en espacio sensor y con 200 surrogates. Aumentar el número de surrogates puede estabilizar p-valores extremos, pero no debe usarse para forzar significación. La siguiente lectura científicamente relevante requiere CSD activo, configuración congelada y comparación sistemática entre sujetos, grupos, sesiones y condiciones.

\subsection{Síntesis técnica de la ejecución}

\begin{ideaclave}
La ejecución de NeuroMIND ya produce artefactos adecuados para el informe: QC, señal filtrada, PSD, potencia por banda, matrices wPLI y contraste inferencial con surrogates/FDR. La lectura más estable es metodológica: el pipeline genera salidas trazables y permite conectar teoría, implementación y resultados. En esta ejecución individual no hay aristas wPLI significativas tras FDR, por lo que la lectura científica queda abierta hasta relanzar el análisis completo con CSD y comparación entre sujetos/condiciones.
\end{ideaclave}

\begin{table}[H]
\centering
\small
\caption{Estado de incorporación de artefactos NeuroMIND al informe. La selección prioriza figuras interpretables en papel y conserva los detalles granulares en la carpeta de resultados.}
\label{tab:results-artifact-selection}
\begin{adjustbox}{max width=0.98\textwidth}
\begin{tabular}{@{}p{0.30\textwidth}p{0.25\textwidth}p{0.36\textwidth}@{}}
\toprule
\textbf{Familia} & \textbf{Uso en informe} & \textbf{Motivo} \\
\midrule
Señal filtrada, PSD, potencia & Incluida como figura & Control visual y lectura espectral directa \\
Topomaps ICA seleccionados & Incluidos como panel QC & Componentes con mayor evidencia automática de artefacto \\
wPLI por bandas & Incluido como figuras & Alfa/theta como foco principal y atlas complementario por bandas \\
Matrices p/q y máscaras & Resumidas en tabla & Mejor lectura inferencial agregada \\
Series ICA completas & Reservadas & Demasiado granulares para esta versión del informe \\
CSV/JSON principales & Citados y copiados parcialmente & Trazabilidad metodológica \\
\bottomrule
\end{tabular}
\end{adjustbox}
\end{table}

Las figuras descartadas para esta versión son las matrices de p/q por banda completas, las máscaras significativas completas y las series temporales ICA individuales. No se excluyen por baja calidad, sino por requerir una explicación demasiado granular para esta sección; quedan preservadas en NeuroMIND y los resúmenes surrogate principales se copiaron a \path{assets/figures/results/} para trazabilidad.
